% Body-only section for the parent's REVTeX report.
% Requires amsmath, amssymb; no theorem package or custom macros are required.
% Bibliography keys: TranquilliSandu2014, LiuMoursiVanderwerff2022.

\section{Mathematical audit and constructive repair}
\label{sec:theory-audit}

The audit separates a mathematical premise from the numerical evidence used to
instantiate it.  Exact carrier coordinates, polynomial resolvents, and finite
causal stage propagation are valid constructions under their declared
structure.  The present evidence does not establish a fundamental failure of
these ideas.  It does identify obligations that an allocation optimization
cannot discharge: the structure must hold on the relevant domain; a residual
must be transported through an inverse or flow-sensitivity bound; and the
result must refer to the original physical norm and requested outputs.
Conversely, temporary storage before an unchanged linear solve and continued
execution of a comparator whose cost has already lost are implementation or
scheduling costs.  Removing them does not require changing the integration
formula or its error threshold.

The distinction is relevant to the interpretation of the existing code.
The coupled solver explicitly labels its original-space residual as
operational rather than certified.  The causal-majorant implementation
contains finite triangular propagation and radius closure.  The conditioned
polynomial path binds a declared polynomial identity, coefficient defects,
and an original-operator residual.  These safeguards should be retained.
Classical Rosenbrock order conditions do not license arbitrary approximate
stage solves or an arbitrary replacement Jacobian; Rosenbrock--W and
Rosenbrock--K methods account for these approximations through their own order
theory.  This distinction is explicit in Sec.~1, following Eqs.~(2a)--(2b), of
Ref.~\cite{TranquilliSandu2014}.

\subsection{A partial-cost stopping theorem for the adapter}
\label{subsec:cost-theorem}

Fix an accepted checkpoint, a final time, an output contract, and an accuracy
qualification.  Both candidates use the same original state and cost units.
Let a completed explicit candidate satisfy this qualification at cost at most
$E$.  Let $L_k$ be the nondecreasing cost already spent by an implicit
comparator after its $k$th checked unit.  Every completion of that same
comparator must have a total cost $I\geq L_k$.  The native success flag supplies
the registered local-tolerance and output qualification; it does not by itself
supply a rigorous physical global-error bound.  If the stronger qualification
is required, it must be established separately.

\paragraph{Theorem 1 (prefix-cost dominance).}
For $0<\rho<1$, the condition
\begin{equation}
 E\leq \rho L_k
 \label{eq:prefix-dominance}
\end{equation}
implies $E\leq\rho I$ for every completion of the implicit comparator.  Hence
the explicit result may be selected, and the comparator may be stopped,
without computing its endpoint.  If the comparator does not complete, it
supplies no competing qualified result; the explicit result retains its
previously established qualification.

\paragraph{Proof.}
The monotonicity premise gives
$E\leq\rho L_k\leq\rho I$ for every comparator completion.  No state from the
incomplete comparator enters this inequality or the explicit result.
Numerical admissibility of the latter was established before the cost test.
\hfill$\square$

The latest baseline already uses the lower-bound argument for an incomplete
comparator.  The implementation refinement is to stop proactively once
Eq.~\eqref{eq:prefix-dominance} becomes decisive, use conservative scalar
arithmetic, and record evidence distinct from a completed same-window
comparison.  A callback-only lower bound is sufficient when all omitted
factorization, solve, and bookkeeping costs are nonnegative.  A count of
callbacks cannot be substituted for elapsed seconds in an inequality about
elapsed seconds.

In floating-point arithmetic, use an upper enclosure of $E$ and a lower
enclosure of $\rho L_k$.  With serial checking after a callback of maximum
weight $U$, a prescribed callback threshold $B$ permits charged work up to
$B+U$, together with any separately charged structural verification.  This is
a bound in the registered callback units, not a wall-clock bound.
An observed fatal callback remains terminal even when an explicit candidate
is available.  A typed cost stop must not conceal a fatal error.  Callback
replay must satisfy the caller's purity or transaction contract, and partial
implicit states are never committed.  The theorem compares the present
qualified observations; it does not guarantee future optimality or a fixed
fraction saved in the total solve, since candidate construction and discarded
work have already been paid.

\subsection{Uniformly convergent reuse of a dissipative shifted system}
\label{subsec:cayley-theorem}

Let $H$ be a finite-dimensional real or complex Hilbert space with a fixed
positive-definite metric $M$.  All norms and adjoints in this subsection refer
to this same metric.  A linear operator $J$ is assumed dissipative:
\begin{equation}
 \operatorname{Re}\langle v,Jv\rangle_M\leq0
 \quad\hbox{for every }v\in H.
 \label{eq:dissipativity}
\end{equation}
Negative real parts of eigenvalues alone are not this hypothesis.  Nor can a
state-dependent error metric be changed silently between the cached and
current systems.

Use the native scaled stage convention
\begin{equation}
 \begin{gathered}
 P=a_0^{-1}I-J,\qquad W=a^{-1}I-J,\\
 a_0>0,\quad a=\gamma h>0.
 \end{gathered}
 \label{eq:native-shifts}
\end{equation}
The shifts $a_0,a$ have units of time; $J,P,W$ have units of inverse time.
Set
\begin{equation}
 \omega_*={2a\over a+a_0},\qquad
 q_*={|a-a_0|\over a+a_0},\qquad
 R_0=(I-a_0J)^{-1}.
 \label{eq:optimal-relaxation}
\end{equation}

\paragraph{Theorem 2 (minimax scalar relaxation).}
Under Eq.~\eqref{eq:dissipativity}, both $P$ and $W$ are invertible.  The
iteration
\begin{equation}
 x_{k+1}=x_k+\omega_*P^{-1}(b-Wx_k)
 \label{eq:relaxed-richardson}
\end{equation}
contracts the error in the $M$ norm by at most $q_*<1$ for every positive shift
ratio.  Moreover, $\omega_*$ minimizes the uniform worst-case norm bound over
real scalar relaxations and all dissipative operators in this class.

\paragraph{Proof.}
If $Wv=0$, taking the real part of the inner product with $v$ gives
$a^{-1}\|v\|_M^2=\operatorname{Re}\langle v,Jv\rangle_M\leq0$, hence $v=0$.
Finite dimensionality implies invertibility.  The same reasoning applies to
$P$.

Write $v=(I-a_0J)u$.  Dissipativity gives
\begin{align}
 \|u+a_0Ju\|_M^2-\|u-a_0Ju\|_M^2
 &=4a_0\operatorname{Re}\langle u,Ju\rangle_M\nonumber\\
 &\leq0.
\end{align}
Therefore the reflected resolvent
\begin{equation}
 C=2R_0-I=(I+a_0J)(I-a_0J)^{-1}
\end{equation}
has norm at most one.  Put $s=a_0/a$.  Direct algebra yields
\begin{align}
 P^{-1}W
 &=I+(s-1)R_0\nonumber\\
 &={1+s\over2}I+{s-1\over2}C,
\end{align}
so the iteration matrix is exactly
\begin{equation}
 T_*=I-\omega_*P^{-1}W
 ={a-a_0\over a+a_0}C.
 \label{eq:cayley-iteration-identity}
\end{equation}
For $e_k=x_k-W^{-1}b$, one has $e_{k+1}=T_*e_k$, proving the contraction.

For minimax sharpness, the class contains $J=0$, for which the norm of the
iteration matrix with scalar relaxation $\omega$ is $|1-\omega s|$.  For the
family $J=-\lambda I$, its limit as $\lambda\to\infty$ is $|1-\omega|$.
Every uniform bound is therefore at least
\begin{equation}
 \max\{|1-\omega s|,|1-\omega|\}.
\end{equation}
Its minimum over real $\omega$ balances the two endpoint errors with opposite
signs, giving $\omega=2/(1+s)$ and value $|1-s|/(1+s)$.  Equation
\eqref{eq:cayley-iteration-identity} attains that uniform bound.
\hfill$\square$

The contraction of a reflected resolvent is a standard monotone-operator
fact, not a claimed novelty of this project.  Fact~2.5 and Eq.~(14) of
Ref.~\cite{LiuMoursiVanderwerff2022} state the corresponding equivalence; in
that convention the monotone operator is $A=-a_0J$.  The contribution here is
the explicit connection to the solver's scaled stage system and a bounded,
testable reuse candidate.  The existing lagged driver already has
Richardson, GMRES, and several staleness criteria.  Theorem~2 supplies an
additional metric-qualified candidate; it does not establish that this
candidate is faster than those existing choices.

\paragraph{Jacobian drift.}
Suppose the cached operator obeys
$\operatorname{Re}\langle v,J_0v\rangle_M\leq-\alpha\|v\|_M^2$ with
$\alpha\geq0$, while the current operator is $J=J_0+\Delta$.
Coercivity and Cauchy--Schwarz give
\begin{equation}
 \|P^{-1}\|_M\leq{a_0\over1+a_0\alpha}.
\end{equation}
Since the perturbed iteration matrix is
$T=T_*+\omega_*P^{-1}\Delta$,
\begin{equation}
 \|T\|_M\leq q_*+
 \omega_*{a_0\|\Delta\|_M\over1+a_0\alpha}=:q.
 \label{eq:cayley-drift}
\end{equation}
This derivation does not assume that $J_0$ and $\Delta$ commute.  A contraction
claim requires $q<1$ and a valid drift bound over the declared domain.

If the implemented inverse action is $\widehat R$, rather than exact
$P^{-1}$, and
$\|\widehat R-P^{-1}\|_M\leq\varepsilon_R$, the iteration also contains this
action error.  If the represented relaxation differs from the mathematical
value by at most $\varepsilon_\omega$, a sufficient additional allowance is
\begin{equation}
 \varepsilon_\omega\|P^{-1}W\|_M+
 |\widehat\omega|\varepsilon_R\|W\|_M.
 \label{eq:inverse-action-allowance}
\end{equation}
For fixed $J$, $\|P^{-1}W\|_M\leq\max(1,a_0/a)$.
A favorable observed residual ratio is not itself a full-space proof of
$\varepsilon_R$ or $q$.  Alternatively, use the iteration heuristically and
certify its final result by an independently valid original-target residual
bound.

\paragraph{Native residual units.}
If the actual $J$ has dissipation margin $\alpha'\geq0$, then for the original
scaled residual $r=b-Wx$,
\begin{equation}
 \|x-W^{-1}b\|_M
 \leq{a\over1+a\alpha'}\|r\|_M.
 \label{eq:scaled-residual-bound}
\end{equation}
Indeed, $\|W^{-1}\|_M\leq a/(1+a\alpha')$ follows from the same coercivity
argument.  Here $r$ has units of state per time.  The factor $a$ must not be
omitted or mixed with the normalized equation using $I-aJ$, whose residual
has state units.  The cached dissipation margin $\alpha$ in
Eq.~\eqref{eq:cayley-drift} and the actual-operator margin $\alpha'$ in
Eq.~\eqref{eq:scaled-residual-bound} are separate hypotheses.

\paragraph{Computational form and scope of verification.}
For $a\leq a_0$, compute $r=a/a_0$, then
\begin{equation}
 \omega_*={2r\over1+r},\qquad q_*={1-r\over1+r}.
\end{equation}
For $a>a_0$, compute $r=a_0/a$ and use $\omega_*=2/(1+r)$ with the same
expression for $q_*$.  This avoids overflow in $a+a_0$ and avoids a ratio
greater than one.  The research Rust primitive encloses the ratio, applies
the monotone formulas with outward scalar arithmetic, includes the
represented-$\omega$ allowance and declared drift, and refuses stale
model/metric/domain identities.  Its scalar argument assumes correctly
rounded basic binary64 operations, gradual underflow, and no
fast-math reassociation.  Tags prevent an identity mismatch; they do not
prove the truth of a caller declaration.  The primitive does not certify an
external floating-point LU action or residual callback, and it is not wired
into a production solver default.  Compilation with Rust 1.94.1 and all
seven targeted primitive tests succeeded.

The Python prototype used actual cached LU factors on 27 finite shifted
systems: diagonal dissipative, damped nonnormal dissipative, and skew-adjoint
families, with shift ratios from $0.1$ to $10$.  All 27 numerical
operator-norm and matrix-identity checks satisfied their diagnostic
tolerances.  The separate residual threshold $10^{-10}$ was met by
\emph{25 of 27} relaxed runs.  Two ratio-10 cases stalled at floating
residual floors and remain recorded as failures; the threshold was not
relaxed.  The largest observed forward error over all 27 cases was
$2.4181321575052657\times10^{-11}$, and the largest matrix-identity discrepancy
was $4.719879393761364\times10^{-13}$.  Exact rational $2\times2$ Wolfram
checks separately confirmed Eq.~\eqref{eq:cayley-iteration-identity}, the
noncommuting drift identity, and the counterexample below.

Unrelaxed iteration diverged in the diagonal ratio-$0.1$ and ratio-$0.25$
cases, while the relaxed candidate converged.  Nevertheless, the uniform
minimax rule generally increased iteration counts in the other families.
The actual spectrum can be easier than the worst-case dissipative class.
For example, for a very stiff negative scalar operator,
$P^{-1}W$ approaches $I$.  Unrelaxed Richardson then has an iteration
matrix approaching zero, whereas a nonunit minimax relaxation leaves
$I-\omega_*I$.  Balancing the zero and infinitely stiff endpoints of an
entire operator class can therefore sacrifice speed on a particular stiff
operator.  If the fixed metric is not the physical Euclidean metric, the
physical norm conversion must also be included; the simple bound acquires
a factor at most $\sqrt{\kappa_2(M)}$ in its initial-error prefactor.
Thus these results establish a larger guaranteed convergence region and an
executable candidate; they establish no universal speedup.  A direct
refactorization or the existing unrelaxed iteration may be cheaper.  The
candidate must remain conditional and be selected using its total measured
cost at the same original-state accuracy.

\subsection{A causal stage bound that includes inverse-action defects}
\label{subsec:inexact-causal-theorem}

Consider finite causal stage equations
\begin{equation}
 W_i u_i=b_i(u_0,\ldots,u_{i-1}),\qquad i=0,\ldots,s-1,
\end{equation}
and candidate stages $\widehat u_i$.  All $W_i$ and approximate inverse
actions $R_i$ are fixed square linear operators on the same finite-dimensional
space.  In a fixed common norm, assume
\begin{align}
 \|I-R_iW_i\|&\leq\delta_i<1,\nonumber\\
 \|R_i r_i\|&\leq d_i,\nonumber\\
 \|R_i[b_i(v)-b_i(w)]\|
 &\leq\sum_{j<i}K_{ij}\|v_j-w_j\|,
 \label{eq:inexact-stage-premises}
\end{align}
where $r_i=b_i(\widehat u_{<i})-W_i\widehat u_i$ and $K_{ij}\geq0$.
The domain of the last inequality must contain the candidate and exact
stage states.  A validated radius-closure argument can establish this
premise; a few samples cannot.

\paragraph{Theorem 3 (defect-aware causal propagation).}
The recurrence
\begin{equation}
 E_i={d_i+\sum_{j<i}K_{ij}E_j\over1-\delta_i}
 \label{eq:inexact-causal-majorant}
\end{equation}
encloses every stage error:
$\|\widehat u_i-u_i\|\leq E_i$.  For the reconstructed output
$y^+=y+\sum_i\beta_i u_i$, its added algebraic error is bounded by
\begin{equation}
 \|\widehat y^+-y^+\|
 \leq\sum_i|\beta_i|E_i+\varepsilon_{\rm reconstruction}.
 \label{eq:stage-output-bound}
\end{equation}
There is no requirement that the full interstage coupling matrix be a
contraction.

\paragraph{Proof.}
Let $D_i=I-R_iW_i$ and $e_i=\widehat u_i-u_i$.  Subtracting exact and
candidate stage equations and multiplying by $R_i$ gives
\begin{equation}
 (I-D_i)e_i
 =-R_i r_i+R_i[b_i(\widehat u_{<i})-b_i(u_{<i})].
\end{equation}
The Neumann series gives $\|(I-D_i)^{-1}\|\leq(1-\delta_i)^{-1}$.
Consequently,
\begin{equation}
 \|e_i\|\leq
 {d_i+\sum_{j<i}K_{ij}\|e_j\|\over1-\delta_i}.
\end{equation}
Induction in causal stage order proves Eq.~\eqref{eq:inexact-causal-majorant}.
The output estimate follows by the triangle inequality.
\hfill$\square$

This formulation extends the existing finite triangular majorant by first
absorbing each diagonal inverse-action defect.  It permits a polynomial
inverse or cached factor without treating that action as exact.  An
unavailable $\delta_i$ leaves an operational estimate, rather than a
certified stage bound.  Full-space operator bounds, truthful domain
declarations, and outward arithmetic are still required when the word
``certified'' is used.

For connection to time integration, suppose the ideal order-$p$ method
satisfies a stability bound
$\operatorname{Lip}(\Phi_h)\leq1+Lh$ and a uniform ideal local defect
$Ch^{p+1}$.  If Eq.~\eqref{eq:stage-output-bound} is at most $Dh^{p+1}$,
then, at the computed time sequence,
\begin{equation}
 e_{n+1}\leq(1+Lh_n)e_n+(C+D)h_n^{p+1}.
\end{equation}
Expanding this recurrence, using
$\prod_n(1+Lh_n)\leq e^{LT}$ and
$\sum_n h_n^{p+1}\leq T h_{\max}^p$, yields
\begin{equation}
 e_N\leq e^{LT}\bigl[e_0+(C+D)T h_{\max}^p\bigr].
\end{equation}
Thus the extra algebraic error preserves order $p$ under these premises.
Stiffness-uniform order requires stiffness-uniform stability and defect
constants; it does not follow from this recurrence alone.  An embedded pair
is not a proof of those constants.

\paragraph{Nonnormal counterexample.}
Set
\begin{equation}
 J=\begin{pmatrix}-1&K\\0&-1\end{pmatrix},\quad
 a=1,\quad b=\begin{pmatrix}0\\\epsilon\end{pmatrix},\quad
 \widehat x=0.
\end{equation}
Both eigenvalues of $J$ are $-1$, and the residual in
$(I-J)x=b$ has norm $\epsilon$.  Direct inversion gives
\begin{equation}
 x=\begin{pmatrix}K\epsilon/4\\\epsilon/2\end{pmatrix}.
\end{equation}
At $K=10^8$, $\epsilon=10^{-8}$, the residual is $10^{-8}$ but the forward
error is approximately $1/4$.  Exact rational evaluation gives
$(1/4,1/200000000)^T$.  The example refutes the inference from a small
residual and negative eigenvalues to a small physical forward error.  It
does not refute the solver or the defect-aware construction.

\subsection{Transport of a structural-chart residual to physical space}
\label{subsec:chart-theorem}

Let
\begin{equation}
 y=S(t)z+c(t)
\end{equation}
with $S$ invertible and continuously differentiable throughout the interval,
and $c$ continuously differentiable.  Define the exact transformed field
\begin{equation}
 g(t,z)=S^{-1}\bigl[f(t,Sz+c)-\dot S z-\dot c\bigr].
 \label{eq:exact-chart-field}
\end{equation}

\paragraph{Theorem 4 (physical residual transport).}
For a differentiable numerical curve $\widehat z$ and its exact
reconstruction $\widehat y=S\widehat z+c$, define
$d_z=\widehat z'-g(t,\widehat z)$.  Then
\begin{equation}
 d_y:=\widehat y'-f(t,\widehat y)=S(t)d_z.
 \label{eq:chart-residual}
\end{equation}
If the implemented transformed field is $\widetilde g$, define instead
$\widetilde d_z=\widehat z'-\widetilde g(t,\widehat z)$ and
$\eta=\widetilde g-g$.  The corresponding identity is
\begin{equation}
 d_y=S(t)(\widetilde d_z+\eta).
 \label{eq:approximate-chart-residual}
\end{equation}

\paragraph{Proof.}
Differentiate the reconstruction:
$\widehat y'=\dot S\widehat z+S\widehat z'+\dot c$.
Subtract $f(t,\widehat y)$ and insert Eq.~\eqref{eq:exact-chart-field}.
This yields Eq.~\eqref{eq:chart-residual}.  Substituting
$\widehat z'-g=\widetilde d_z+\eta$ yields
Eq.~\eqref{eq:approximate-chart-residual}.
\hfill$\square$

To obtain a global physical error bound, assume that both exact and numerical
curves lie in a domain on which the vector field has a valid one-sided
Lipschitz bound $\ell(t)$ in the fixed physical Hilbert norm.  For
$e=\widehat y-y$, the differential norm inequality is
\begin{equation}
 D^+\|e\|\leq\ell(t)\|e\|+\|d_y\|.
\end{equation}
With $\Gamma(T,t)=\exp(\int_t^T\ell(s)\,ds)$, the comparison lemma gives
\begin{align}
 \|e(T)\|\leq{}&\Gamma(T,t_0)\|e(t_0)\|\nonumber\\
 &+\int_{t_0}^{T}\Gamma(T,t)\|S(t)\|
       \bigl(\|\widetilde d_z\|+\|\eta\|\bigr)\,dt.
 \label{eq:chart-global-bound}
\end{align}
Piecewise reconstructions require the corresponding propagated jump terms.
An approximate reconstruction also needs its own defect allowance; the
identity for exact $S$ does not certify a floating implementation of $S$.

The implication for the declared chirp and cofactor paths is precise.
The analytic construction can remove a known carrier, but a small
transformed residual is insufficient when the lift is large in the physical
norm.  Whole-domain structural accuracy, phase evaluation, reconstruction,
and flow bounds remain provider obligations.  Endpoint samples cannot replace
these hypotheses.  The existing declared-family implementations are not
evidence for arbitrary noncommuting carriers or arbitrary black-box vector
fields.  Those extensions require a derived time-ordered transformation and
a bound on its defect, rather than a change in storage or loop scheduling.

\subsection{Theory-to-production admission}

The four results have different implementation status.  Prefix-cost
dominance justifies the bounded comparator optimization while preserving the
existing numerical qualification.  The dissipative shifted-system result
has a proved conditional convergence statement, exact finite CAS checks, a
tested scalar Rust primitive, and finite numerical experiments; its iteration
regressions and two residual-target failures exclude unconditional promotion.
The defect-aware causal recurrence reuses the existing majorant structure
but still requires original-domain inverse and coupling bounds.  The chart
identity specifies how physical error must be transported; it does not supply
the missing whole-domain bounds for a particular physical client.

Accordingly, a production integration of the new relaxation candidate must
retain the current target equation, bind a declared dissipative metric and
drift domain, account for numerical inverse-action and residual-evaluation
errors, and compare against current Richardson, GMRES, and refactorization
paths at the same original-state accuracy.  The exact allocation optimizations
and comparator scheduling changes are separate advances.  Their measured
benefit must not be used to claim that an untested structural or theoretical
candidate has already become faster.
